\chapter{Primitive Roots and the Multiplicative Group}\label{ch:orders}

\begin{center}\itshape
The decimal expansion of $\frac17=0.\overline{142857}$ repeats every $6$ digits, and $\frac1{17}$ every $16$.\\
Why do some primes $p$ give the longest possible period $p-1$, and are there infinitely many of them? Nobody knows.
\end{center}


The decimal digits of $1/7$ repeat in a block of length $6$. The reason has nothing to do with long division being lucky. Modulo $7$, successive powers of $10$ give
\[
3,\quad2,\quad6,\quad4,\quad5,\quad1.
\]
Every nonzero residue occurs before the sequence returns to $1$. The smallest return time is called the \emph{order} of $10$ modulo $7$. Here it is $6$, the largest it could be.

The story changes when the modulus changes. Modulo $8$, each odd residue squares to $1$. Thus no unit can have powers visiting all four units: every orbit has at most two terms. What holds in a prime field need not hold in an arbitrary ring of remainders. That small contrast motivates the central theorem of the chapter: the nonzero residue classes modulo a prime form a cyclic multiplicative group.

Three tasks must be kept separate. First, show that some element visits all units. Secondly, find such an element for a particular modulus. Thirdly, find the exponent that reaches a chosen unit. The first is an existence theorem, the second a search, and the third a discrete logarithm problem. A proof of existence does not automatically give a fast algorithm for the other two.

The order of an element is our first exact measure of periodicity. Later, characters will use the same cycles to record residue classes, and quadratic residues will be the even powers inside them.


\begin{reminder}[A change of coordinate]
The additive group of remainders is generated by $1$. This chapter asks when
multiplication admits an equally simple coordinate: a unit $g$ whose powers
visit all the other units. The answer depends on the modulus. For a prime it is
always yes; modulo $8$ it is already no. We will separate existence, finding a
generator, and computing the exponent of a given unit. They are three different
problems, with three very different costs.
\end{reminder}

\section{Orders}

% t1-pass2 motive 953a9dc1
Euler's theorem says $a^{\varphi(n)}\equiv 1$, but the least such positive exponent is usually smaller. That least exponent is the order. It divides every other exponent that works, and it is the lattice of the cyclic subgroup generated by $a$. Volume~6's characters are homomorphisms out of this group; a primitive root is a generator.

\begin{definition}
Let $\gcd(a,n)=1$. The \emph{order} $\operatorname{ord}_n(a)$ is the least $k\ge1$ with $a^k\equiv1\pmod n$.
\end{definition}

\begin{proposition}\label{prop:order}
Let $k=\operatorname{ord}_n(a)$. Then $a^m\equiv1\pmod n$ if and only if $k\mid m$. In particular $k\mid\varphi(n)$. Moreover $\operatorname{ord}_n(a^j)=\dfrac k{\gcd(k,j)}$.
\end{proposition}
\begin{proof}
Because $\gcd(a,n)=1$, the class of $a$ is a unit,
so its powers are defined for every integer exponent.
The set
\[
H=\{m\in\Z:a^m\equiv1\pmod n\}
\]
contains $0$ and is closed under subtraction. It is
therefore a subgroup of $(\Z,+)$.
Euler's theorem supplies the positive element
$\varphi(n)\in H$. Let $k$ be the least positive
element of $H$, the order of $a$.
The subgroup theorem from Chapter~\ref{ch:divisibility}
gives $H=k\Z$, so $a^m\equiv1$ precisely when $k\mid m$.
In particular $k\mid\varphi(n)$.

Now let $j$ be an integer. The order of $a^j$ is
the least positive $m$ satisfying $k\mid jm$.
Write $d=\gcd(k,j)$, $k=dk'$ and $j=dj'$.
Since $\gcd(k',j')=1$, Gauss's lemma gives
$k'\mid m$ whenever $k'\mid j'm$.
Conversely, every multiple of $k'$ satisfies
$k\mid jm$. Therefore the least such $m$ is
\[
\operatorname{ord}_n(a^j)=k'=\frac{k}{\gcd(k,j)}.
\]
The argument includes $j=0$, when the power is $1$
and its order is $1$.
\end{proof}

The definition makes sense because Euler's theorem supplies at least one
positive exponent returning to $1$. Negative powers in the proof are powers
of the multiplicative inverse, which exists precisely because $\gcd(a,n)=1$.
For $j=0$, the order formula gives $k/\gcd(k,0)=1$, as it should: $a^0=1$.

\paragraph{The first return controls every return.}
Modulo $13$, the powers of $3$ repeat as
\[
1,3,9,1,3,9,\dots,
\]
so its order is $3$.
If $3^{100}\equiv3\pmod{13}$, this is not a coincidence:
$100=3\cdot33+1$. In general, division $m=qk+r$, $0\le r<k$, gives
$a^m\equiv a^r$; the least-return property forces $r=0$ exactly when the
power is $1$. This is an alternative to the subgroup argument above.

Here is a useful distinction. Euler says $a^{\varphi(n)}\equiv1$; he does
not say that $\varphi(n)$ is the \emph{first} such exponent. Modulo $8$,
$3^2\equiv1$ although $\varphi(8)=4$. An order is a period, not a deadline.

A \emph{primitive root} modulo $n$ is an $a$ with $\operatorname{ord}_n(a)=\varphi(n)$, that is, a generator of the group $(\Z/n\Z)^\times$, which is then cyclic.

% t1-pass2 examples 953a9dc1
\begin{examplex}[the order of $2$ modulo $7$]
The powers of $2$ modulo $7$ are $2,4,1$. The order is $3$, which divides $\varphi(7)=6$ and is not $6$. So $2$ is not a primitive root modulo $7$. The powers of $3$ are $3,2,6,4,5,1$, and the order is $6$.
\end{examplex}

\begin{examplex}[orders divide]
Since $2^3\equiv 1\pmod 7$, also $2^6\equiv 1$ and $2^9\equiv 1$, but $2^1$ and $2^2$ are not $1$. The positive exponents that work are exactly the multiples of $3$. This is the content of the first proposition of the section, read on a single residue.
\end{examplex}

\section{\texorpdfstring{$(\Z/p\Z)^\times$}{(Z/pZ)*} is cyclic}

% t1-pass2 motive 57059d76
A finite subgroup of the multiplicative group of a field is cyclic, and the proof is a counting of elements of each order. For the field $\Z/p\Z$ this says there is a primitive root. Figure~\ref{fig:t1-cycle} is the generator $3$ modulo $7$. Volume~3 proves the same counting statement for every finite field, and Volume~13 uses a primitive root to build the cyclotomic extension.

\begin{think}
List the powers of $2$ modulo $13$. Do they reach every non-zero class? Then compute the order of each of $1,\dots,12$ modulo $13$ and count how many elements have order $1,2,3,4,6,12$. Compare these counts with $\varphi(d)$. Can you guess the general statement, and see why a polynomial of degree $d$ having at most $d$ roots might be the key?
\end{think}

\begin{theorem}[Gauss]\label{thm:cyclic}
For $p$ prime, $(\Z/p\Z)^\times$ is cyclic: there are exactly $\varphi(p-1)$ primitive roots modulo $p$.
\end{theorem}
\begin{proof}
For $d\mid p-1$, let $N(d)$ be the number of elements of order $d$. Suppose $N(d)>0$, and let $a$ have order $d$. Then $1,a,\dots,a^{d-1}$ are $d$ distinct roots of $X^d-1$ in the field $\Z/p\Z$, and a polynomial of degree $d$ has at most $d$ roots, so these are all the elements $x$ with $x^d=1$. Those of order exactly $d$ are the $a^j$ with $\gcd(j,d)=1$ (Proposition~\ref{prop:order}): $N(d)=\varphi(d)$.
So $N(d)\le\varphi(d)$ for every $d\mid p-1$. But every element has some order dividing $p-1$, so $\sum_{d\mid p-1}N(d)=p-1=\sum_{d\mid p-1}\varphi(d)$ (Proposition~\ref{prop:phisum}). Hence $N(d)=\varphi(d)$ for all $d$, in particular $N(p-1)=\varphi(p-1)\ge1$.
\end{proof}

% t1-pass2 why
\paragraph{Why this works.}
The polynomial $X^d-1$ has at most $d$ roots modulo a prime, so the number of elements whose order divides $d$ is at most the sum of $\varphi$ over the divisors of $d$, with equality in the count only when some element has order exactly $d$. Taking $d=p-1$ produces a primitive root.

\paragraph{Why the field hypothesis matters.}
The polynomial root bound used in the proof follows by repeated division by
$X-a$: if $a$ is a root, write $P=(X-a)Q$; any other root $b$ satisfies
$(b-a)Q(b)=0$, and since $b-a\ne0$ is invertible in a field, $Q(b)=0$.
Induction on the degree gives the bound. Over $\Z/8\Z$, the polynomial
$X^2-1$ has the four roots $1,3,5,7$: a degree-two polynomial has staged a
small rebellion. The failure is exactly that non-zero factors can have product
zero in this ring. Thus the counting proof is exploiting a field property,
not merely a clever arrangement of residues.

There is also a logical point in the argument: we did \emph{not} assume that
an element of order $d$ exists for every divisor $d$. We first proved that
$N(d)$ is either $0$ or $\varphi(d)$. The equality of the two total sums then
forces every deficit to vanish. In particular the largest order exists.
For $p=2$ the same argument includes the one-element group, with
$\varphi(1)=1$ and primitive root $1$.

\begin{figure}[ht]
\centering
\fig{ch05-orders}
\caption{Orders of units modulo $13$. For each divisor $d$ of $12$, the bar
height is $N(d)=\varphi(d)$: $1,1,2,2,2,4$. The twelve units are partitioned
by their \emph{exact} orders; merely knowing $x^d=1$ counts units whose order
divides $d$, which is a different question. The rightmost bar gives the four
primitive roots.}
\label{fig:ordercounts}
\end{figure}

Figure~\ref{fig:primroot} draws the theorem for $p=13$: once the units are placed around a circle in the order of the powers of a primitive root, multiplication becomes rotation.

\begin{figure}[ht]
\centering
\fig{ch05-primroot}
\caption{The twelve units modulo $13$, placed in the order $2^0,2^1,\dots,2^{11}$. Left: multiplying by the primitive root $2$ moves one step and visits every unit. Right: $3=2^4$, so multiplying by $3$ jumps four steps, and the units split into four cycles of length $\operatorname{ord}_{13}(3)=3$. The shaded units $2,6,11,7$ are the $\varphi(12)=4$ primitive roots, the powers $2^j$ with $\gcd(j,12)=1$.}
\label{fig:primroot}
\end{figure}

\paragraph{Finding a generator without an exhausting roll call.}
Let $g$ be a unit modulo $n$ and write $M=\varphi(n)$. Then $g$ has order
$M$ if and only if
\[
 g^{M/q}\not\equiv1\pmod n
 \quad\text{for every prime }q\mid M.
\]
Indeed, if its order $d$ is a proper divisor of $M$, the integer $M/d>1$
has a prime divisor $q$, and then $d\mid M/q$. Conversely any failed test
shows that the order divides a smaller exponent. Notice that this criterion
requires the prime divisors of $M$, not a list of all its divisors.

For $n=13$, $M=12$ has prime divisors $2$ and $3$. We find
$2^6\equiv-1$ and $2^4\equiv3$ modulo $13$, so $2$ is primitive.
By contrast $4^6\equiv1$, so $4$ fails immediately; in fact its order is $6$.
This is a rigorous certificate of success or failure, and it uses just two
modular exponentiations. The hint to Exercise~\ref{ex:findroots} is now an
algorithm we have proved.

% t1-pass2 examples 57059d76
\begin{examplex}[a generator modulo $7$]
The residue $3$ has order $6=\varphi(7)$. Its powers visit every nonzero class, in the order $1,3,2,6,4,5$ if one starts at $3^0$. Multiplication by $3$ is a single $6$-cycle. The residue $2$ generates only the subgroup $\{1,2,4\}$.
\end{examplex}

\begin{figure}[ht]
\centering
\fig{ch05-cycle}
\caption{Multiplication by $3$ modulo $7$ runs through every nonzero residue and returns to $1$ after six steps. This is a picture of a primitive root.}
\label{fig:t1-cycle}
\end{figure}

\begin{examplex}[how many generators]
Once $3$ is known to be primitive modulo $7$, the other primitive roots are $3^k$ with $\gcd(k,6)=1$, namely $k=1$ and $k=5$. Modulo $7$ one has $3^2\equiv 2$, $3^3\equiv -1$, $3^4\equiv 4$ and $3^5\equiv 5$. The primitive roots are $3$ and $5$, and there are $\varphi(6)=2$ of them.
\end{examplex}

\section{Which \texorpdfstring{$(\Z/n\Z)^\times$}{(Z/nZ)*} are cyclic?}

% t1-pass2 motive 5e93c0f8
The group of units modulo $n$ is cyclic only for a short list of $n$. The obstruction is visible: too many solutions of $x^2=1$. The complementary structure theorem, for powers of $2$, says the group is a product of two cyclic groups generated by $-1$ and $5$. Volume~13's local square theorem at the prime $2$ is the statement that the squares are exactly one of the four cosets modulo $8$.

\begin{think}
Square each of the units $1,3,5,7$ modulo $8$. What do you get? Why does this show that $(\Z/8\Z)^\times$ has no generator? Which other moduli might suffer from the same problem?
\end{think}

\begin{lemma}\label{lem:lift}
Let $p$ be an odd prime and $g$ an integer with $g^{p-1}=1+pt$, $p\nmid t$. Then for every $k\ge1$, $g^{(p-1)p^{k-1}}=1+p^kt_k$ with $p\nmid t_k$.
\end{lemma}
\begin{proof}
Induction on $k$. If $x=1+p^kt_k$ with $p\nmid t_k$, then
$x^p=1+p^{k+1}t_k+\binom p2p^{2k}t_k^2+\dots$. Every further term is divisible by $p^{k+2}$: $\binom p2p^{2k}$ because $p\mid\binom p2$ ($p$ odd) and $2k+1\ge k+2$, and the others because $jk\ge k+2$ for $j\ge3$. So $x^p=1+p^{k+1}t_{k+1}$ with $t_{k+1}\equiv t_k\not\equiv0\pmod p$.
\end{proof}

\begin{theorem}\label{thm:cyclicn}
$(\Z/n\Z)^\times$ is cyclic if and only if $n\in\{1,2,4\}$ or $n=p^k$ or $n=2p^k$ with $p$ an odd prime and $k\ge1$.
\end{theorem}
\begin{proof}
\emph{$n=p^k$.} Let $g$ be a primitive root modulo $p$. If $g^{p-1}\equiv1\pmod{p^2}$, replace $g$ by $g+p$, still primitive modulo $p$: by the binomial theorem $(g+p)^{p-1}\equiv g^{p-1}+(p-1)pg^{p-2}\equiv1-pg^{p-2}\not\equiv1\pmod{p^2}$. So we may assume $g^{p-1}=1+pt$ with $p\nmid t$. Let $m=\operatorname{ord}_{p^k}(g)$. It divides $\varphi(p^k)=(p-1)p^{k-1}$, and is a multiple of $\operatorname{ord}_p(g)=p-1$, so $m=(p-1)p^j$ with $j\le k-1$. If $j<k-1$, then $g^{(p-1)p^{k-2}}\equiv1\pmod{p^k}$, contradicting the lemma (applied with $k-1$). So $m=\varphi(p^k)$.

\emph{$n=2p^k$.} By the Chinese remainder theorem $(\Z/2p^k\Z)^\times\cong(\Z/2\Z)^\times\times(\Z/p^k\Z)^\times$, and the first factor is trivial. The cases $n=1,2,4$ are checked by hand ($3$ generates modulo $4$); at $n=1$ the unit group is understood as the one-element group.

\emph{The others are not cyclic.} A cyclic group of even order $m$, generated by $g$, has exactly two solutions of $x^2=1$: $(g^j)^2=1$ iff $m\mid2j$ iff $j\in\{0,\frac m2\}$ modulo $m$. For $n>2$, $\varphi(n)$ is even, so it suffices to exhibit four square roots of $1$. If $8\mid n$: $\pm1$ and $2^{a-1}\pm1$ modulo $2^a$ ($a\ge3$) give four, which combine with $1$ modulo the odd part by the Chinese remainder theorem. If $n$ has two distinct odd prime factors, or $n=4p^k$, then $n=ab$ with $a,b>2$ coprime, and the pairs $(\pm1,\pm1)$ give four.
\end{proof}

% t1-pass2 why
\paragraph{Why this works.}
A cyclic group of even order has exactly two solutions of $x^2=1$. Exhibiting four solutions, by the Chinese remainder theorem when $n$ has two odd prime factors or when $8$ divides $n$, shows the group cannot be cyclic. The positive cases are a primitive root modulo $p$, lifted to $p^k$ by a binomial expansion that raises the valuation of $g^{\varphi(p^k)/p}-1$ by exactly one at each step.

\paragraph{Watching a primitive root lift.}
Take $p=3$ and $g=2$. Since $2^2=1+3$, the lemma applies with $t=1$.
The orders modulo $3,9,27$ are therefore $2,6,18$:
\[
 2^2=1+3,\qquad 2^6=64=1+9\cdot7,
 \qquad 2^{18}\equiv1+27\pmod{81}.
\]
The first non-zero power of $3$ in $g^{2\cdot3^{k-1}}-1$ increases by
exactly one each time we cube. It cannot skip a level. This is what prevents
the order from stopping short of $2\cdot3^{k-1}$ modulo $3^k$.
If the first candidate has $g^{p-1}\equiv1\pmod{p^2}$, changing it to
$g+p$ repairs precisely that first step; its class modulo $p$ does not change.

The non-cyclic cases are worth comparing directly. Modulo $15$, the four
square roots of $1$ are $1,4,11,14$, the Chinese remainder combinations of
$\pm1$ modulo $3$ and $5$. Modulo $8$ they are $1,3,5,7$.
A cyclic group of even order has room for just one non-identity element whose
square is $1$. These four roots are an obstruction one can see, without
trying every possible generator.

% t1-pass2 examples 5e93c0f8
\begin{examplex}[four square roots of $1$ modulo $8$]
The units modulo $8$ are $1,3,5,7$, and each squares to $1$. A cyclic group of even order has only two solutions of $x^2=1$. Hence $(\Z/8\Z)^\times$ is not cyclic. The same four roots survive modulo every higher power of $2$, which is why those groups are not cyclic either.
\end{examplex}

\begin{examplex}[a primitive root modulo $9$]
Modulo $9$, $\varphi(9)=6$. The powers of $2$ are $2,4,8,7,5,1$. The order is $6$, so $2$ is a primitive root modulo $9$, in agreement with the theorem for odd prime powers. The residue $4=2^2$ has order $3$ and is not primitive.
\end{examplex}

\section{Discrete logarithms}

% t1-pass2 motive 2b385d7a
A primitive root turns multiplication into addition of exponents. Solving $x^k\equiv a$ becomes a linear congruence in the exponent. Computing the power is easy and recovering the exponent is not, which is the asymmetry cryptography uses; the exercises here stay with moduli small enough that the exponent can be found by hand.

\begin{think}
Find $i$ with $2^i\equiv5\pmod{13}$. How did you do it? Now imagine $p$ has $600$ digits: how long would your method take, compared with computing $2^i\bmod p$ from $i$?
\end{think}

If $g$ is a primitive root modulo $p$, every unit is $g^i$ for a unique $i\in\Z/(p-1)\Z$, its \emph{index} or \emph{discrete logarithm} $\operatorname{ind}_g$. It turns multiplication into addition: $\operatorname{ind}(ab)=\operatorname{ind}a+\operatorname{ind}b\pmod{p-1}$, so multiplicative questions become linear ones modulo $p-1$ (Exercise~\ref{ex:kth}).
Computing $g^i$ from $i$ is fast (square and multiply); recovering $i$ from $g^i$ is believed to be hard in appropriately chosen large groups. A large $p$ alone is not enough: the factorisation of $p-1$ and the chosen subgroup matter, as does the algorithm available. This asymmetry is the basis of the Diffie--Hellman key exchange.

\paragraph{Solving a power equation in the exponent coordinate.}
For $x^4\equiv3\pmod{13}$, use $g=2$, so $3=2^4$ and write $x=2^u$.
We need $4u\equiv4\pmod{12}$, or $u\equiv1\pmod3$. The four exponents
$1,4,7,10$ give the four roots
\[
 x\equiv2,3,11,10\pmod{13}.
\]
Each fourth power is $3$; these are all the roots because every non-zero
class has a unique exponent modulo $12$. The exponent lives modulo $p-1$,
while the original residue lives modulo $p$. Confusing those two clocks is
the standard trap in discrete-logarithm calculations.

\begin{think}
Would the equation $x^4\equiv2\pmod{13}$ have a solution? The exponent
congruence would be $4u\equiv1\pmod{12}$, impossible because its left side
is divisible by $4$. Equivalently, $2^3\not\equiv1\pmod{13}$.
Now try the same question for $x^3\equiv5\pmod{13}$, using $5=2^9$.
The exponent condition is $u\equiv3\pmod4$, giving $x\equiv8,11,7$.
\end{think}

% t1-pass2 examples 2b385d7a
\begin{examplex}[a logarithm modulo $13$]
The residue $2$ is a primitive root modulo $13$, since $2^6=64\equiv -1$ and $2^4=3\not\equiv 1$. Now $2^4\equiv 3$, so the discrete logarithm of $3$ base $2$ is $4$. Then $3^3=27\equiv 1\pmod{13}$? $27-2\cdot 13=1$, yes, but the order of $3$ is $3$, not $12$. The logarithm is an exponent, not a claim that $3$ itself is primitive.
\end{examplex}

\begin{examplex}[a power equation]
Solve $x^4\equiv 3\pmod{13}$. With $x=2^u$ and $3=2^4$, the congruence becomes $4u\equiv 4\pmod{12}$, so $u\equiv 1\pmod 3$. The solutions are $u=1,4,7,10$, hence $x\equiv 2,3,11,10$. There are $\gcd(4,12)=4$ solutions, as the general count predicts.
\end{examplex}

\section{Units modulo a power of two}

% t1-pass2 motive 485346e0
The group $(\Z/2^k\Z)^\times$ is the one cyclic case the odd primes do not explain. It is generated by $-1$ and by $5$, and the squares are exactly the units congruent to $1$ modulo $8$. That single congruence is the finite reason a $2$-adic unit is a square in $\Z_2$ precisely under the same condition.

Theorem~\ref{thm:cyclicn} shows that $(\Z/2^k\Z)^\times$ is not cyclic for $k\ge3$, by exhibiting four solutions of $x^2\equiv1$. The group is still completely explicit. It is generated by $-1$ and by $5$, and those two generators commute and overlap only at $1$.

\begin{theorem}\label{thm:twopower}
Let $k\ge3$. Every unit modulo $2^k$ can be written uniquely as
\[
(-1)^a5^b\pmod{2^k},
\]
with $a\in\{0,1\}$ and $0\le b<2^{k-2}$.
\end{theorem}
\begin{proof}
There are $\varphi(2^k)=2^{k-1}$ units. The order of $5$ is found from the $2$-adic valuation of its powers. For every integer $m\ge0$,
\[
v_2\bigl(5^{2^m}-1\bigr)=m+2.
\]
The case $m=0$ is $5-1=4$. If the claim holds for $m$, write $5^{2^m}=1+2^{m+2}t$ with $t$ odd. Then
\[
5^{2^{m+1}}=(1+2^{m+2}t)^2=1+2^{m+3}t+2^{2m+4}t^2.
\]
The last term is divisible by $2^{m+4}$, while $t$ odd forces the middle term to have valuation exactly $m+3$. The valuation of the whole difference from $1$ is therefore $m+3$.

In particular $v_2(5^{2^{k-2}}-1)=k$, so $5^{2^{k-2}}\equiv1\pmod{2^k}$, and $v_2(5^{2^{k-3}}-1)=k-1$, so $5^{2^{k-3}}\not\equiv1\pmod{2^k}$. The order of $5$ divides $\varphi(2^k)=2^{k-1}$ and is a power of $2$. It divides $2^{k-2}$ and does not divide $2^{k-3}$, hence equals $2^{k-2}$. The powers $5^0,\dots,5^{2^{k-2}-1}$ are distinct. None of them is congruent to $-1$, because every power of $5$ is congruent to $1$ modulo $4$, while $-1\equiv3\pmod4$. The $2^{k-1}$ classes $(-1)^a5^b$ are therefore distinct, and they exhaust the units.
\end{proof}

% t1-pass2 why
\paragraph{Why this works.}
The valuation $v_2(5^{2^m}-1)=m+2$ says that the order of $5$ modulo $2^k$ is exactly $2^{k-2}$. Powers of $5$ are $1$ modulo $4$, so none of them is $-1$. Counting then matches $\varphi(2^k)$, and every unit appears exactly once.

\begin{corollary}\label{cor:mod8squares}
Let $k\ge3$. A unit modulo $2^k$ is a square if and only if it is congruent to $1$ modulo $8$.
\end{corollary}
\begin{proof}
By Theorem~\ref{thm:twopower} a unit is $(-1)^a5^b$. Its square is $5^{2b}$, an even power of $5$, and every even power arises in this way. Modulo $8$ one has $5^1\equiv5$ and $5^2\equiv1$, so $5^b\equiv1\pmod8$ precisely when $b$ is even, while $-5^b\equiv3$ or $7\pmod8$. Thus the unit is a square if and only if $a=0$ and $b$ is even, which is exactly the condition of being congruent to $1$ modulo $8$.
\end{proof}

The four odd classes modulo $8$ are the four cosets of the subgroup of squares. This is why a $2$-adic square is detected modulo $8$ and not modulo a higher power of $2$: once a unit is a square modulo $8$, the corollary says it is already a square modulo every $2^k$. Volume~13 then promotes that root to a square in the $2$-adic integers, where the derivative condition of Hensel's lemma can be checked at the first admissible precision. The finite group explains the modulus; the completion explains the remaining digits.

\begin{examplex}[the units modulo $16$]
Here $2^{4-2}=4$, so the powers of $5$ are $1,5,9,13$, and the units are
\[
\begin{aligned}
1&=5^0,&5&=5^1,&9&=5^2,&13&=5^3,\\
15&=-1,&11&=-5,&7&=-9,&3&=-13.
\end{aligned}
\]
The squares are the even powers $1$ and $9$, which are exactly the units congruent to $1$ modulo $8$. For instance $7\equiv7\pmod8$ is not a square, while $17\equiv1\pmod8$ must be a square modulo every $2^k$; modulo $32$ one square root is $7$, since $49-17=32$.
\end{examplex}

% t1-pass2 examples 485346e0
\begin{examplex}[the units modulo $16$]
The powers of $5$ modulo $16$ are $1,5,9,13$. Multiplying by $-1$ fills in $15,11,7,3$. These are all eight units. The squares are $1$ and $9$, the two classes congruent to $1$ modulo $8$. In particular $7$ is not a square modulo $16$, and neither is $7$ modulo $8$.
\end{examplex}

\begin{examplex}[the order of $5$ modulo $32$]
Here $2^{5-2}=8$. One checks $5^2=25$, $5^4\equiv 625\equiv 17\pmod{32}$ and $5^8=(1+16)^2=1+32+256\equiv 1\pmod{32}$, while $5^4=17\not\equiv 1$. The order is exactly $8$. This is the case $k=5$ of the valuation computation $v_2(5^{2^m}-1)=m+2$.
\end{examplex}

\begin{beauty}
Where this leads. Primitive roots are the reason $(\Z/p\Z)^\times$ has $\varphi(p-1)$ generators, which is the counting step in the proof that there are $\varphi(m)$ Dirichlet characters in Chapter~\ref{ch:characters}. The structure modulo $2^k$ is the input to the $2$-adic square theorem of Volume~13 and to the local class field theory of the same volume at the prime $2$. Artin's conjecture on primitive roots, still open, is a question about how often this generator exists for a fixed base as $p$ varies; Volume~5 can prove such statements only under a Riemann hypothesis.
\end{beauty}

\subsection*{Proof workshop: decide an order by testing divisors}

Compute the order of $2$ modulo $7$. The powers are
$2,4,1$, so the order is $3$. Fermat's theorem only tells
us that $2^6\equiv1$; it does not claim that $6$ is the
smallest possible exponent.

For a unit $a$ whose order is $r$, a power $a^k$ equals
$1$ exactly when $r$ divides $k$: divide $k$ by $r$, use
$a^r=1$, and observe that a smaller positive remainder
cannot give $1$. Thus the order divides the size of a
finite unit group. For example, modulo $7$ every order
divides $6$. To check whether an element has order $6$,
it is enough to test that the powers with exponents $3$
and $2$ are not $1$. For $3$, these are $3^3\equiv6$
and $3^2\equiv2$; both differ from $1$. So $3$ has
order $6$. This is a practical test, not a guess from the
first few digits.


\section{Exercises}

\begin{exercise}[$\star$]\label{ex:findroots}
Find a primitive root modulo $7$, $11$ and $13$, and list all primitive roots modulo $13$.
\end{exercise}

\begin{exercise}[$\star$]\label{ex:period}
Let $p\ne2,5$ be prime. Show that the decimal expansion of $\frac1p$ is periodic with least period $\operatorname{ord}_p(10)$. Check the periods of $\frac17$ and $\frac1{17}$.
\end{exercise}

\begin{exercise}[$\star\star$]\label{ex:fermatfactor}
For $n\ge0$, show that every prime factor $p$ of the Fermat number $2^{2^n}+1$ satisfies $p\equiv1\pmod{2^{n+1}}$. (This is how Euler found $641\mid2^{32}+1$.)
\end{exercise}

\begin{exercise}[$\star\star$]\label{ex:powersum}
Let $p$ be prime and $k\ge0$. Show that $\displaystyle\sum_{x=1}^{p-1}x^k\equiv\begin{cases}-1&\text{if }(p-1)\mid k,\\0&\text{otherwise}\end{cases}\pmod p$.
\end{exercise}

\begin{exercise}[$\star\star$]\label{ex:kth}
Let $p$ be prime, $k\ge1$ an integer, $p\nmid a$, and $d=\gcd(k,p-1)$. Show that $x^k\equiv a\pmod p$ has a solution if and only if $a^{(p-1)/d}\equiv1\pmod p$, and that it then has exactly $d$ solutions.
\end{exercise}

\begin{exercise}[$\star\star\star$]\label{ex:gausswilson}
\emph{Gauss' generalisation of Wilson's theorem.} Let $n>2$. Show that the product of all units modulo $n$ is $\equiv-1\pmod n$ if $(\Z/n\Z)^\times$ is cyclic, and $\equiv1$ otherwise. \emph{Hint: pair each unit with its inverse, and study the subgroup $\{x:x^2=1\}$.}
\end{exercise}

\begin{exercise}[$\star\star\star$]\label{ex:cyclotomic}
\emph{Infinitely many primes $\equiv1\pmod n$.} For a positive integer $n$, define the cyclotomic polynomials by $X^n-1=\prod_{d\mid n}\Phi_d(X)$, that is, $\Phi_n=\prod_{\gcd(k,n)=1}\left(X-e^{2ik\pi/n}\right)$.
\begin{enumerate}[label=(\alph*)]
\item Show by induction that $\Phi_n$ is monic with integer coefficients, and that $\Phi_n(0)=\pm1$ for $n\ge2$.
\item Let $p$ be a prime with $p\nmid n$ and $p\mid\Phi_n(a)$. Show that $\operatorname{ord}_p(a)=n$, hence $p\equiv1\pmod n$. \emph{Hint: otherwise $a$ would be a double root of $X^n-1$ modulo $p$.}
\item Deduce that there are infinitely many primes $p\equiv1\pmod n$.
\end{enumerate}
\end{exercise}

\begin{exercise}[$\star$]\label{ex:t1-units16}
Using Theorem~\ref{thm:twopower}, write every unit modulo $16$ as $(-1)^a5^b$ and decide which of them are squares.
\end{exercise}

\begin{exercise}[$\star\star$]\label{ex:t1-square-cosets}
Let $k\ge3$. Show that the subgroup of squares in $(\Z/2^k\Z)^\times$ has index $4$, and that its four cosets are precisely the sets of units congruent to $1,3,5,7$ modulo $8$.
\end{exercise}

\begin{exercise}[$\star$]\label{ex:t1-gen7}
Show that $3$ is a primitive root modulo $7$ and that $2$ is not. List the other primitive root.
\end{exercise}

\begin{beauty}
Theorem~\ref{thm:cyclic} is the reason finite fields are so pleasant: their multiplicative groups are cyclic, a fact Volume~3 proves for every finite field $\mathbb F_q$ by the same counting argument. \emph{Artin's conjecture} (1927) predicts that $10$ is a primitive root modulo a proportion $\approx37.4\%$ of the primes. More generally every integer other than $-1$ and perfect squares should be primitive for a positive proportion of primes, but that proportion can depend on the integer: $37.4\%$ is not a universal answer. For a specified base such as $10$ the assertion remains open unconditionally; Hooley proved it under the generalised Riemann hypothesis for the relevant number fields (1967). Exercise~\ref{ex:cyclotomic} is the special case $a\equiv1$ of Dirichlet's theorem on primes in progressions, proved in full in Volume~6.
\end{beauty}

\begin{exercise}[$\star\star$]\label{pb:smith}
(Smith, 1875.)
\begin{enumerate}[label=(\alph*)]
\item Show that $\det\bigl(\gcd(i,j)\bigr)_{1\le i,j\le n}=\varphi(1)\varphi(2)\cdots\varphi(n)$.
\item More generally, let $f$ be any function on the positive integers and $g=f*\mathbf1$, that is $g(m)=\sum_{d\mid m}f(d)$.
Show that $\det\bigl(g(\gcd(i,j))\bigr)_{1\le i,j\le n}=f(1)\cdots f(n)$, and compute $\det\bigl(\gcd(i,j)^2\bigr)_{1\le i,j\le n}$.
\end{enumerate}
\end{exercise}

\begin{exercise}[$\star\star$]\label{pb:lambert}
(Lambert series.) Let $(a_n)_{n\ge1}$ be complex numbers with $|a_n|\le Cn^K$ for some constants $C,K$.
\begin{enumerate}[label=(\alph*)]
\item Show that for $|x|<1$, $\displaystyle\sum_{n\ge1}a_n\frac{x^n}{1-x^n}=\sum_{m\ge1}\Bigl(\sum_{d\mid m}a_d\Bigr)x^m$.
\item Deduce the sums of $\sum_n\varphi(n)\frac{x^n}{1-x^n}$, $\sum_n\mu(n)\frac{x^n}{1-x^n}$ and $\sum_n\lambda(n)\frac{x^n}{1-x^n}$,
where $\lambda(n)=(-1)^{\Omega(n)}$ is Liouville's function ($\Omega(n)$ counts the prime factors of $n$ with multiplicity).
\item Show that $\displaystyle\sum_{n\ge1}\frac{\varphi(n)}{2^n-1}=2$ and $\displaystyle\sum_{n\ge1}\frac{\mu(n)}{2^n-1}=\frac12$.
\end{enumerate}
\end{exercise}

\begin{exercise}[$\star\star$]\label{pb:ramsum}
For integers $q\ge1$ and $n$, the \emph{Ramanujan sum} is $c_q(n)=\sum e^{2i\pi an/q}$, over $1\le a\le q$ with
$\gcd(a,q)=1$.
\begin{enumerate}[label=(\alph*)]
\item Show that $c_q(n)=\sum_{d\mid\gcd(q,n)}\mu(q/d)\,d$. In particular $c_q(n)$ is an integer.
\item Show Hölder's formula: with $g=\gcd(q,n)$,
\[
c_q(n)=\mu\Bigl(\frac qg\Bigr)\frac{\varphi(q)}{\varphi(q/g)}.
\]
\item Let $q$ and $r$ divide $N$. Show that $\frac1N\sum_{n=1}^Nc_q(n)c_r(n)$ is $\varphi(q)$ if $q=r$ and $0$ otherwise.
\end{enumerate}
\end{exercise}

\begin{exercise}[$\star\star$]\label{pb:phipm1}
\begin{enumerate}[label=(\alph*)]
\item Show that $\Phi_n(1)=e^{\Lambda(n)}$ for $n\ge2$: it is $p$ if $n$ is a power of a prime $p$, and $1$ otherwise.
\item Compute $\Phi_n(-1)$ for every $n\ge3$.
\item Let $n\ge2$ not be a prime power and $\zeta=e^{2i\pi/n}$. Show that $1-\zeta$ is invertible in the ring $\Z[\zeta]$.
\end{enumerate}
\end{exercise}

\begin{exercise}[$\star\star$]\label{pb:gcdmean}
\emph{The law of the gcd.} Let $a,b$ be independent and uniform in $\{1,\dots,N\}$.
\begin{enumerate}[label=(\alph*)]
\item Show that $P(\gcd(a,b)=g)\to\frac6{\pi^2g^2}$ for every $g\ge1$.
\item Show that $\mathbb E[\gcd(a,b)]=\frac6{\pi^2}\log N+O(1)$.
\item Show that $\mathbb E[\log\gcd(a,b)]\to-\frac{\zeta'(2)}{\zeta(2)}=0.5699\ldots$
\item Show that $\mathbb E[\operatorname{lcm}(a,b)]\sim\frac{\zeta(3)}{4\zeta(2)}N^2=0.1827\ldots N^2$.
\end{enumerate}
\end{exercise}

\section{Hints}

\begin{hint}{ex:findroots}
$g$ is a primitive root modulo $p$ if and only if $g^{(p-1)/q}\not\equiv1$ for every prime $q\mid p-1$. Once you have one primitive root $g$, the others are the $g^j$ with $\gcd(j,p-1)=1$.
\end{hint}
\begin{hint}{ex:period}
In the long division of $1$ by $p$, the remainder after $j$ steps is $10^j\bmod p$. When does it come back to $1$?
\end{hint}
\begin{hint}{ex:fermatfactor}
From $2^{2^n}\equiv-1\pmod p$, find the exact order of $2$ modulo $p$, then use that it divides $p-1$.
\end{hint}
\begin{hint}{ex:powersum}
Write the units as $g^0,\dots,g^{p-2}$ for a primitive root $g$. The sum becomes geometric with ratio $g^k$.
\end{hint}
\begin{hint}{ex:kth}
Take discrete logarithms: $x=g^\xi$, $a=g^\alpha$. The equation becomes $k\xi\equiv\alpha\pmod{p-1}$, a linear congruence from Chapter~\ref{ch:congruences}.
\end{hint}
\begin{hint}{ex:gausswilson}
After pairing each unit with its inverse, only the solutions of $x^2=1$ remain. How many are there when the group is cyclic, and what is their product when there are $2^r\ge4$ of them?
\end{hint}
\begin{hint}{ex:cyclotomic}
(a) Divide $X^n-1$ by $\prod_{d\mid n,\,d<n}\Phi_d$, a monic polynomial with integer coefficients. (b) If $\operatorname{ord}_p(a)=d<n$, then modulo $p$ the class of $a$ is a root of $X^d-1$, hence of some $\Phi_e$ with $e\mid d$, and also of $\Phi_n$: a double root of $X^n-1$. Look at the derivative. (c) If $p_1,\dots,p_r$ were all the primes $\equiv1\pmod n$, evaluate $\Phi_n$ at a large multiple of $m=np_1\cdots p_r$: the value is $\equiv\pm1\pmod m$.
\end{hint}
\begin{hint}{ex:t1-units16}
The order of $5$ modulo $16$ is $4$. Square the four powers and compare them with the units congruent to $1$ modulo $8$.
\end{hint}
\begin{hint}{ex:t1-square-cosets}
Corollary~\ref{cor:mod8squares} names the squares. The odd residue classes modulo $8$ partition the units, and each class contains the same number of them.
\end{hint}

\begin{hint}{ex:t1-gen7}
The order divides $6$, so it is enough to rule out orders $1$, $2$ and $3$. The other generator is $3^5$ modulo $7$.
\end{hint}

\begin{hint}{pb:smith}
$\gcd(i,j)=\sum_{d\mid i,\ d\mid j}\varphi(d)$. Write the matrix as $ADA^{\mathsf T}$ with $A$ triangular.
\end{hint}

\begin{hint}{pb:lambert}
Expand $\frac{x^n}{1-x^n}=\sum_{k\ge1}x^{kn}$, justify the regrouping, and recall the Dirichlet convolutions $\varphi*\mathbf1$,
$\mu*\mathbf1$ and $\lambda*\mathbf1$.
\end{hint}

\begin{hint}{pb:ramsum}
(a) Sort the $a\in\{1,\dots,q\}$ by $\gcd(a,q)$ and use Möbius inversion. (b) Both sides are multiplicative in $q$; check
prime powers. (c) Expand both sums and use that $a/q$ and $b/r$, in lowest terms, are equal modulo $1$ only if they are
the same fraction.
\end{hint}

\begin{hint}{pb:phipm1}
(a) Evaluate $\frac{x^n-1}{x-1}=\prod_{d\mid n,d>1}\Phi_d$ at $1$ and take logarithms. (b) Prove $\Phi_{2m}(x)=\Phi_m(-x)$
for odd $m>1$, and $\Phi_{pm}(x)=\Phi_m(x^p)$ when $p\mid m$. (c) $\Phi_n(1)$ is a product of factors $1-\zeta^a$.
\end{hint}

\begin{hint}{pb:gcdmean}
(b) $\gcd(a,b)=\sum_{d\mid a,\ d\mid b}\varphi(d)$. (d) Pull out the gcd and use Möbius for the coprime part.
\end{hint}

\section{Solutions}

\begin{solution}{ex:findroots}
\idea To test whether $g$ is primitive modulo $p$, we need not compute all its powers: by Proposition~\ref{prop:order} its order divides $p-1$, so it is primitive iff $g^{(p-1)/q}\not\equiv1$ for each prime $q\mid p-1$. Try $g=2,3,\dots$. Once one primitive root is known, all of them are its powers with exponent prime to $p-1$.
\writeup Modulo $7$: powers of $3$ are $3,2,6,4,5,1$, so $3$ is primitive. Modulo $11$: $2$ has order dividing $10$, and $2^2=4$, $2^5=32\equiv-1$, so the order is $10$. Modulo $13$: $2^4=16\equiv3$, $2^6=64\equiv-1$, so $\operatorname{ord}(2)$ divides $12$ but not $4$ or $6$: it is $12$. The primitive roots modulo $13$ are $2^j$ with $\gcd(j,12)=1$ (Proposition~\ref{prop:order}): $2^1,2^5,2^7,2^{11}\equiv2,6,11,7$. There are $\varphi(12)=4$.
\end{solution}

\begin{solution}{ex:period}
\idea In long division, the digits are determined by the successive remainders, and the remainder after $j$ steps is $10^j\bmod p$. So the expansion repeats exactly when the powers of $10$ modulo $p$ do, i.e.\ with period the order of $10$.
\writeup Long division of $1$ by $p$ produces the remainders $r_j=10^j\bmod p$ (each step multiplies the remainder by $10$ and reduces), and the $(j+1)$-th digit is $\lfloor10r_j/p\rfloor$, a function of $r_j$. Since $10$ is invertible modulo $p$, the sequence $10^j\bmod p$ is purely periodic with least period $\operatorname{ord}_p(10)$, and so are the digits; the least period of the digits is not smaller, because $\frac1p=\frac{A}{10^k-1}$ with $A$ an integer (period $k$) means exactly $p\mid10^k-1$. Modulo $7$, $10\equiv3$ is primitive: period $6$. Modulo $17$: $10^2\equiv-2$, $10^8\equiv16\equiv-1$, so the order is $16$.
\end{solution}

\begin{solution}{ex:fermatfactor}
\idea $p\mid2^{2^n}+1$ says $2^{2^n}\equiv-1$, hence $2^{2^{n+1}}\equiv1$: we know the order of $2$ divides a power of $2$, and $-1\ne1$ pins it down exactly. Then order divides $p-1$.
\writeup $2^{2^n}\equiv-1\pmod p$, so $2^{2^{n+1}}\equiv1$. The order of $2$ divides $2^{n+1}$ but not $2^n$ (else $2^{2^n}\equiv1\not\equiv-1$, as $p$ is odd). So $\operatorname{ord}_p(2)=2^{n+1}$, which divides $p-1$.
\end{solution}

\begin{solution}{ex:powersum}
\idea With a primitive root $g$, the sum over all units becomes a geometric sum $\sum_i(g^k)^i$, which is easy: it depends on whether the ratio $g^k$ is $1$.
\writeup Let $g$ be a primitive root. The $x$ run over $g^0,\dots,g^{p-2}$, so the sum is $\sum_{i=0}^{p-2}(g^k)^i$. If $(p-1)\mid k$, $g^k\equiv1$ and the sum is $p-1\equiv-1$. Otherwise $g^k\not\equiv1$ and the geometric sum is $\frac{g^{k(p-1)}-1}{g^k-1}\equiv0$, since $g^{k(p-1)}\equiv1$ and $g^k-1$ is invertible.
\end{solution}

\begin{solution}{ex:kth}
\idea The discrete logarithm turns the multiplicative equation $x^k=a$ into the linear congruence $k\xi\equiv\alpha\pmod{p-1}$, which Chapter~\ref{ch:congruences} solves completely. Then translate the condition $d\mid\alpha$ back into a condition on $a$.
\writeup Write $a=g^\alpha$ and $x=g^\xi$. The equation becomes $k\xi\equiv\alpha\pmod{p-1}$, solvable if and only if $d\mid\alpha$, with exactly $d$ solutions modulo $p-1$ (Proposition~\ref{prop:lincong}). And $d\mid\alpha$ iff $(p-1)\mid\alpha\frac{p-1}d$ iff $a^{(p-1)/d}=g^{\alpha(p-1)/d}\equiv1$.
\end{solution}

\begin{solution}{ex:gausswilson}
\idea Copy the proof of Wilson's theorem: pairing $x$ with $x^{-1}$ reduces the product to the product of the self-inverse elements, the solutions of $x^2=1$. These form a group where every element has order $\le2$, which is a vector space over $\mathbb F_2$; the sum of all vectors of such a space is $0$ unless it is a line.
\writeup Pairing each unit with its inverse, all pairs $\{x,x^{-1}\}$ with $x\ne x^{-1}$ contribute $1$, so the product equals the product of the elements of $S=\{x:x^2=1\}$. $S$ is a subgroup in which every element is its own inverse. If $|S|=2$, $S=\{1,-1\}$ and the product is $-1$. Otherwise $|S|>2$. Written additively, $S$ is an abelian group in which every element has order at most $2$, that is, a vector space over $\mathbb F_2$, of some dimension $r\ge2$. The sum of all vectors of $\mathbb F_2^r$ is $0$ when $r\ge2$, because each coordinate equals $1$ in exactly $2^{r-1}$ vectors, an even number. So the product is $1$. Finally, by the proof of Theorem~\ref{thm:cyclicn}, $|S|=2$ exactly when $(\Z/n\Z)^\times$ is cyclic.
\end{solution}

\begin{solution}{ex:cyclotomic}
\idea This is Euclid's argument with $\Phi_n$ in place of $X+1$: we need a polynomial whose prime factors (apart from those of $n$) are all $\equiv1\pmod n$. Part (b) gives this: a prime dividing $\Phi_n(a)$ makes $a$ an element of order exactly $n$ modulo $p$, because a smaller order would make $a$ a double root of $X^n-1$, impossible when $p\nmid n$ (the derivative test). Then evaluate $\Phi_n$ at a multiple of $n$ and of all known such primes, so that the new prime factor is new.
\writeup (a) $\Phi_1=X-1$. If the claim holds for all proper divisors of $n$, then $\Phi_n=(X^n-1)/\prod_{d\mid n,d<n}\Phi_d$ is the quotient of a monic integer polynomial by a monic integer polynomial, and the division algorithm stays in $\Z[X]$ (each step divides by the leading coefficient $1$). The remainder is zero because the factorisation over $\C$ already shows exact divisibility; uniqueness of quotient and remainder identifies the integer quotient with $\Phi_n$. Evaluating at $0$: $-1=\prod_{d\mid n}\Phi_d(0)=-\prod_{d\mid n,d>1}\Phi_d(0)$, and by induction all factors except $\Phi_n(0)$ are $\pm1$, so $\Phi_n(0)=\pm1$.

(b) $a^n\equiv1\pmod p$ since $\Phi_n(a)\mid a^n-1$. Suppose $d=\operatorname{ord}_p(a)<n$. Then $a$ is a root modulo $p$ of $X^d-1=\prod_{e\mid d}\Phi_e$, hence of some $\Phi_e$ with $e\mid d$, $e<n$, $e\mid n$ ($\Z/p\Z$ is a field). So $a$ is a root of both $\Phi_e$ and $\Phi_n$, two distinct factors of $X^n-1$: it is a double root of $X^n-1$ modulo $p$, so it is also a root of the derivative $nX^{n-1}$ (differentiate $X^n-1=(X-a)^2Q$ in $(\Z/p\Z)[X]$). But $p\nmid n$ and $p\nmid a$: contradiction. So $\operatorname{ord}_p(a)=n$, and $n\mid p-1$.

(c) For $n=1$, the conclusion is Euclid's theorem. Let $n\ge2$ and suppose $p_1,\dots,p_r$ are all the primes $\equiv1\pmod n$. Let $m=np_1\cdots p_r$. For a large integer $t$, $|\Phi_n(mt)|>1$ (a non-constant polynomial tends to infinity), so it has a prime factor $q$. Since $\Phi_n(mt)\equiv\Phi_n(0)=\pm1\pmod m$, $q$ divides neither $n$ nor any $p_i$, and $q\mid\Phi_n(mt)$. By (b), $q\equiv1\pmod n$, a new prime: contradiction.
\end{solution}

\begin{solution}{ex:t1-units16}
\idea Theorem~\ref{thm:twopower} gives eight explicit representatives, and Corollary~\ref{cor:mod8squares} says that squaring is the same test as reduction modulo $8$.
\writeup Modulo $16$ the powers of $5$ are $1$, $5$, $9$ and $13$.
The eight units are $1$, $15$, $5$, $11$, $9$, $7$, $13$ and $3$.
Their squares are $1$ and $9=5^2$, since $5^2\equiv9$, $5^4\equiv1$, and the classes with a minus sign square to the same powers. Both squares are congruent to $1$ modulo $8$, and no other unit in the list is.
\end{solution}

\begin{solution}{ex:t1-square-cosets}
\idea The corollary already names the squares as one of the four odd classes modulo $8$. Counting how the units fall into those classes gives the index.
\writeup By Corollary~\ref{cor:mod8squares} the squares are exactly the units congruent to $1$ modulo $8$. In a complete set of residues modulo $2^k$ each class modulo $8$ contains $2^{k-3}$ integers, and the four odd classes consist entirely of units. The units congruent to $1$ modulo $8$ are closed under multiplication and inverses, so they form a subgroup of order $2^{k-3}$ and of index $4$. The other three odd classes are its cosets.
\end{solution}

\begin{solution}{ex:t1-gen7}
\idea Compute powers until the residue $1$ appears, and compare the first return with $p-1$.
\writeup Modulo $7$, the powers of $2$ are $2,4,1$, so the order is $3$, not $6$. The powers of $3$ are $3,2,6,4,5,1$, so the order is $6$. Then $3^5\equiv 5$, and $\gcd(5,6)=1$, so $5$ is the other primitive root. There are $\varphi(6)=2$ of them.
\end{solution}

\begin{solution}{pb:smith}
\idea A matrix whose entries are a function of $\gcd(i,j)$ is a sum over common divisors, which is a matrix product: one
factor records ``$d$ divides $i$'', the other ``$d$ divides $j$''. The divisibility matrix is triangular with ones on the
diagonal, so the determinant is the product of the diagonal weights.

\writeup (b) Let $A$ be the $n\times n$ matrix with $A_{id}=1$ if $d\mid i$ and $0$ otherwise, and $D=\operatorname{diag}
(f(1),\dots,f(n))$. Then
\[
(ADA^{\mathsf T})_{ij}=\sum_{d=1}^nA_{id}f(d)A_{jd}=\sum_{d\mid i,\ d\mid j}f(d)=\sum_{d\mid\gcd(i,j)}f(d)=g(\gcd(i,j)),
\]
since the common divisors of $i$ and $j$ are the divisors of $\gcd(i,j)$, all at most $n$. The matrix $A$ is lower triangular
($d\mid i$ implies $d\le i$) with $A_{ii}=1$, so $\det A=1$ and $\det\bigl(g(\gcd(i,j))\bigr)=\det D=f(1)\cdots f(n)$ (Volume~1,
Chapter~7).

(a) is the case $f=\varphi$, since $\sum_{d\mid m}\varphi(d)=m$ (Volume~2, Chapter~6).

For $\gcd(i,j)^2$: we need $f$ with $f*\mathbf1=\mathrm{id}^2$, that is $f=\mu*\mathrm{id}^2$ by Möbius inversion, the Jordan
function $J_2(m)=\sum_{d\mid m}\mu(d)\bigl(\frac md\bigr)^2=m^2\prod_{p\mid m}\bigl(1-\frac1{p^2}\bigr)$. So
$\det\bigl(\gcd(i,j)^2\bigr)=\prod_{m=1}^nJ_2(m)$; for $n=4$ this is $1\cdot3\cdot8\cdot12=288$.
\tools Volume~2, Chapter~6 (Dirichlet convolution, Möbius inversion); Volume~1, Chapter~7 (determinants of triangular matrices).
\leads H.~J.~S.~Smith found this in 1875. The same factorisation shows the matrix $(\gcd(i,j))$ is positive definite, and
the asymptotic behaviour of its eigenvalues is still studied (Lindqvist and Seip, 1998).
\end{solution}

\begin{solution}{pb:lambert}
\idea $\frac{x^n}{1-x^n}=x^n+x^{2n}+x^{3n}+\cdots$ has a term $x^m$ exactly when $n\mid m$. So the Lambert series collects, in
front of $x^m$, the sum of $a_d$ over the divisors $d$ of $m$: it is the power series of $a*\mathbf1$. Choosing $a$ with a
known $a*\mathbf1$ gives closed forms.

\writeup (a) For $|x|<1$ the family $\bigl(a_nx^{kn}\bigr)_{n,k\ge1}$ is absolutely summable:
\[
\begin{aligned}
\sum_n|a_n|\sum_k|x|^{kn}
&=\sum_n|a_n|\frac{|x|^n}{1-|x|^n}\\
&\le\frac1{1-|x|}\sum_nCn^K|x|^n<\infty
\end{aligned}.
\]
So it may be summed in any
order (Volume~1, Chapter~3): summing over $k$ first gives the Lambert series, and grouping the terms with $kn=m$ gives
$\sum_m\bigl(\sum_{d\mid m}a_d\bigr)x^m$.

(b) Since $\varphi*\mathbf1=\mathrm{id}$, $\mu*\mathbf1=[m=1]$ (Volume~2, Chapter~6):
\[
\sum_n\varphi(n)\frac{x^n}{1-x^n}=\sum_mmx^m=\frac x{(1-x)^2},\qquad\sum_n\mu(n)\frac{x^n}{1-x^n}=x.
\]
For $\lambda$: it is completely multiplicative, so $\lambda*\mathbf1$ is multiplicative, and on a prime power
$\sum_{j=0}^a\lambda(p^j)=1-1+\dots\pm1$ is $1$ if $a$ is even and $0$ if $a$ is odd. So $(\lambda*\mathbf1)(m)=1$ if $m$ is a
square and $0$ otherwise, and $\sum_n\lambda(n)\frac{x^n}{1-x^n}=\sum_{j\ge1}x^{j^2}$. (Similarly $\Lambda*\mathbf1=\log$ gives
$\sum_m(\log m)x^m$.)

(c) At $x=\frac12$, $\frac{x^n}{1-x^n}=\frac1{2^n-1}$, so the two sums are $\frac{1/2}{1/4}=2$ and $\frac12$.
\tools Volume~2, Chapter~6 (convolution identities); Volume~1, Chapter~3 and Volume~5, Chapter~2 (summable families, power series).
\leads The Lambert series of $\lambda$ is a theta function; Lambert series of divisor functions are Eisenstein series, the
first modular forms (Volume~6, Chapter~12, and Volume~13).
\end{solution}

\begin{solution}{pb:ramsum}
\idea The sum over \emph{all} $a$ is a geometric sum, $q$ or $0$. Sorting the $a$ by $\gcd(a,q)$ writes this full sum as a
sum of Ramanujan sums over the divisors of $q$, and Möbius inversion undoes it. For (b), multiplicativity reduces
everything to prime powers. For (c), think of $c_q(n)$ as a sum over the fractions with denominator exactly $q$.

\writeup Write $e(x)=e^{2i\pi x}$ and $\eta_q(n)=\sum_{a=1}^qe(an/q)$, which is $q$ if $q\mid n$ and $0$ otherwise.

(a) For $a\in\{1,\dots,q\}$ with $\gcd(a,q)=q/e$, write $a=\frac qeb$ with $1\le b\le e$ and $\gcd(b,e)=1$; then
$e(an/q)=e(bn/e)$. This is a bijection between $\{1,\dots,q\}$ and the pairs $(e,b)$ with $e\mid q$, $b$ prime to $e$. So
$\eta_q(n)=\sum_{e\mid q}c_e(n)$, and Möbius inversion (Volume~2, Chapter~6) gives
\[
c_q(n)=\sum_{e\mid q}\mu(q/e)\eta_e(n)=\sum_{e\mid q,\ e\mid n}\mu(q/e)\,e.
\]

(b) For fixed $n$, the function $h(e)=e\cdot\mathbf 1_{e\mid n}$ is multiplicative in $e$, so $c_q(n)=(\mu*h)(q)$ is
multiplicative in $q$ (Volume~2, Chapter~6). The right-hand side of Hölder's formula is multiplicative in $q$ too: if $q=q_1q_2$
with $\gcd(q_1,q_2)=1$, then $g=g_1g_2$ with $g_i=\gcd(q_i,n)$, and $\mu$ and $\varphi$ are multiplicative. So take $q=p^a$,
$a\ge1$, and $g=p^b$ with $b=\min(a,v_p(n))$. In (a) only $e=p^a$ and $e=p^{a-1}$ have $\mu(q/e)\ne0$:
\begin{itemize}
\item if $b=a$, $c_{p^a}(n)=p^a-p^{a-1}=\varphi(p^a)$;
\item if $b=a-1$, only $e=p^{a-1}$ divides $n$, and $c_{p^a}(n)=-p^{a-1}$;
\item if $b\le a-2$, neither divides $n$, and $c_{p^a}(n)=0$.
\end{itemize}
Hölder's formula gives $\mu(1)\varphi(p^a)$, then $\mu(p)\frac{\varphi(p^a)}{\varphi(p)}=-p^{a-1}$, then $0$, since
$\mu(p^{a-b})=0$ for $a-b\ge2$. They agree.

(c) The numbers $c_r(n)$ are integers, hence real, so $c_r(n)=\overline{c_r(n)}=\sum_be(-bn/r)$. Then
\[
\frac1N\sum_{n=1}^Nc_q(n)c_r(n)=\sum_{a,b}\frac1N\sum_{n=1}^Ne\Bigl(n\Bigl(\frac aq-\frac br\Bigr)\Bigr),
\]
with $a$ prime to $q$ and $b$ prime to $r$. Since $q,r\mid N$, $u=e(\frac aq-\frac br)$ is an $N$-th root of unity, and the
inner average is $1$ if $u=1$ and $0$ otherwise. And $u=1$ means $\frac aq\equiv\frac br\pmod1$; both fractions are in
lowest terms in $\left]0,1\right]$, so this means $q=r$ and $a=b$. The total is $\varphi(q)$ if $q=r$, and $0$ otherwise.
\tools Volume~2, Chapter~6 (Möbius inversion, multiplicative functions); Volume~1, Chapter~5 (roots of unity).
\leads Ramanujan (1918) expanded arithmetic functions in series of his sums, like Fourier series:
$\frac{\sigma(n)}n=\frac{\pi^2}6\sum_{q\ge1}\frac{c_q(n)}{q^2}$. Part (c) is the orthogonality
that makes such expansions work. Ramanujan sums also give the main terms in the circle method and in sums of the form
$\sum_n\Lambda(n)\Lambda(n+h)$ (Volume~14).
\end{solution}

\begin{solution}{pb:phipm1}
\idea At $x=1$, the identity $1+x+\dots+x^{n-1}=\prod_{d\mid n,d>1}\Phi_d(x)$ becomes $n=\prod\Phi_d(1)$, a
multiplicative version of $\log n=\sum_{d\mid n}\Lambda(d)$: by Möbius, $\log\Phi_n(1)=\Lambda(n)$. At $x=-1$, two rules
move the index: substituting $-x$ doubles odd indices, and substituting $x^p$ multiplies by $p$ an index that $p$
already divides. Each rule is proved by comparing roots.

\writeup (a) For $n\ge2$, $\frac{x^n-1}{x-1}=\prod_{d\mid n,d>1}\Phi_d(x)$; at $x=1$, $n=\prod_{d\mid n,d>1}\Phi_d(1)$. All values
are positive (the left side is, and by induction each $\Phi_d(1)$ is). So $L(d)=\log\Phi_d(1)$ for $d\ge2$, $L(1)=0$,
satisfies $\sum_{d\mid n}L(d)=\log n=\sum_{d\mid n}\Lambda(d)$ for all $n\ge1$ (Volume~2, Chapter~6). By Möbius inversion, $L=\Lambda$:
$\Phi_n(1)=e^{\Lambda(n)}$.

(b) Two rules. \emph{(i) If $m>1$ is odd, $\Phi_{2m}(x)=\Phi_m(-x)$.} The roots of $\Phi_m(-x)$ are $-\zeta$ for $\zeta$ of order
$m$, and $-\zeta$ has order $2m$ ($m$ odd); there are $\varphi(m)=\varphi(2m)$ of them, distinct, and $\Phi_m(-x)$ is monic
since $\varphi(m)$ is even. Two monic polynomials of the same degree with the same simple roots are equal.
\emph{(ii) If $p\mid m$, $\Phi_{pm}(x)=\Phi_m(x^p)$.} If $z^p$ has order $m$, the order $e$ of $z$ satisfies
$e/\gcd(e,p)=m$, so $e=m$ or $pm$, and $e=m$ would need $\gcd(m,p)=1$. So the $p\varphi(m)$ distinct roots of $\Phi_m(x^p)$ have
order $pm$, and $\varphi(pm)=p\varphi(m)$ when $p\mid m$: same conclusion.

Now let $n\ge3$. If $n$ is odd, $\Phi_n(-1)=\Phi_{2n}(1)=1$ by (i) and (a). If $n=2m$ with $m\ge3$ odd,
$\Phi_n(-1)=\Phi_m(1)$, which is $p$ if $m$ is a power of $p$ and $1$ otherwise. If $4\mid n$, (ii) with $p=2$ gives
$\Phi_n(-1)=\Phi_{n/2}(1)$, which is $2$ if $n$ is a power of $2$ and $1$ otherwise. In short, for $n\ge3$:
$\Phi_n(-1)=p$ if $n=2p^k$ with $p$ an odd prime, $2$ if $n$ is a power of $2$, and $1$ otherwise.

(c) $\Phi_n(1)=\prod_{\gcd(a,n)=1}(1-\zeta^a)=1$ by (a). Each factor is in $\Z[\zeta]$, so $1-\zeta$ times the product of
the other factors is $1$.
\tools Volume~2, Chapters~6 and~4 (Möbius inversion, cyclotomic polynomials); Volume~1, Chapter~6 (polynomials and their roots).
\leads The units $\frac{1-\zeta^a}{1-\zeta}$ and the $1-\zeta$ of (c) are the \emph{cyclotomic units}. Kummer used them in
his attack on Fermat's last theorem, and they are linked to the values $L(1,\chi)$ by the class number formula (Volume~13).
When $n=p^k$, on the contrary, $1-\zeta$ generates the unique prime above $p$.
\end{solution}

\begin{solution}{pb:gcdmean}
\idea $\gcd(a,b)=g$ means $g\mid a$, $g\mid b$ and the quotients are coprime: probability $\frac1{g^2}\cdot\frac6{\pi^2}$. The tail
$\sum\frac1{g^2}\cdot g$ diverges logarithmically, which is why the mean grows like $\log N$, while $\log g$ is still integrable.

\writeup (a) $\gcd(a,b)=g$ if and only if $a=ga'$, $b=gb'$ with $a',b'\le N/g$ coprime. By Exercise~6.8, the number of coprime
pairs in $\{1,\dots,M\}^2$ is $\frac6{\pi^2}M^2+O(M\log M)$, so $P(\gcd=g)=\frac1{N^2}\bigl(\frac6{\pi^2}\frac{N^2}{g^2}+O(\frac Ng\log N)\bigr)\to\frac6{\pi^2g^2}$.

(b) By Gauss's identity $\sum_{d\mid m}\varphi(d)=m$ (Proposition~3.7), $\gcd(a,b)=\sum_{d\mid a,\,d\mid b}\varphi(d)$. So
\[
\sum_{a,b\le N}\gcd(a,b)=\sum_{d\le N}\varphi(d)\Bigl\lfloor\frac Nd\Bigr\rfloor^2=N^2\sum_{d\le N}\frac{\varphi(d)}{d^2}+O\Bigl(N\sum_{d\le N}\frac{\varphi(d)}d\Bigr),
\]
because $\lfloor N/d\rfloor^2=N^2/d^2+O(N/d)$. The error is $O(N^2)$. For the main term, $\Phi(x)=\sum_{d\le x}\frac{\varphi(d)}d=\frac6{\pi^2}x+O(\log x)$
(Volume~6, Exercise~2.1), and Abel summation gives
\[
\sum_{d\le N}\frac{\varphi(d)}{d^2}=\frac{\Phi(N)}N+\int_1^N\frac{\Phi(x)}{x^2}\,dx=\frac6{\pi^2}\log N+O(1).
\]
Divide by $N^2$.

(c) $\mathbb E[\log\gcd]=\sum_{g\le N}\log g\,P(\gcd=g)$. Each term tends to $\frac{6\log g}{\pi^2g^2}$, and is bounded by
$\log g\cdot\frac{\lfloor N/g\rfloor^2}{N^2}\le\frac{\log g}{g^2}$, which is summable. By dominated convergence the expectation tends
to $\frac6{\pi^2}\sum_g\frac{\log g}{g^2}=\frac{-\zeta'(2)}{\zeta(2)}$, since $-\zeta'(2)=\sum_g\frac{\log g}{g^2}=0.93754\ldots$

(d) Since $\operatorname{lcm}(a,b)=\frac{ab}{\gcd(a,b)}$, grouping by $g=\gcd(a,b)$,
\[
\sum_{a,b\le N}\operatorname{lcm}(a,b)=\sum_{g\le N}g\,C\Bigl(\frac Ng\Bigr),\qquad C(M)=\sum_{\substack{a',b'\le M\\\gcd(a',b')=1}}a'b'.
\]
By Möbius, $C(M)=\sum_{d\le M}\mu(d)d^2S(M/d)^2$ with $S(x)=\sum_{n\le x}n=\frac{x^2}2+O(x)$; so
$C(M)=\frac{M^4}4\sum_{d\le M}\frac{\mu(d)}{d^2}+O(M^3\log M)=\frac{6}{\pi^2}\frac{M^4}4+O(M^3\log M)$. Then
\[
\sum_{g\le N}g\,C\Bigl(\frac Ng\Bigr)=\frac6{\pi^2}\frac{N^4}4\sum_{g\le N}\frac1{g^3}+O\Bigl(N^3\log N\sum_g\frac1{g^2}\Bigr)=\frac{\zeta(3)}{4\zeta(2)}N^4+O(N^3\log N).
\]
Divide by $N^2$.

\tools Exercise~6.8 (coprime pairs), Proposition~3.7 (Gauss's identity) and Theorem~6.4 (Möbius inversion); Volume~6, Chapter~2; Volume~4, Chapter~5 (dominated
convergence).
\leads Cesàro studied the mean of the gcd in the 1880s, and the error term in (b) is related to the Dirichlet divisor problem.
The values in (c) and (d) mean that the gcd of two random integers is ``typically'' $e^{0.57}\approx1.77$ but has an infinite
mean in the limit, a heavy tail like the one in Problem~\ref{pb:valuation}(c).
\end{solution}

\section*{Sources and further reading}

\begin{itemize}
\item K.~Ireland and M.~Rosen, \emph{A Classical Introduction to Modern Number Theory}, 2nd ed., Springer, 1990, Chapter~4. The structure of $(\Z/n\Z)^\times$ and power residues.
\item C.~F.~Gauss, \emph{Disquisitiones Arithmeticae} (1801), Section~III; English translation by A.~A.~Clarke, Springer, 1986. Where primitive roots, indices and the generalised Wilson theorem first appear.
\item M.~R.~Murty, Artin's conjecture for primitive roots, \emph{The Mathematical Intelligencer} 10 (1988), 59--67. A readable survey of the open problem quoted at the end of the chapter.
\item C.~Hooley, On Artin's conjecture, \emph{Journal f\"ur die reine und angewandte Mathematik} 225 (1967), 209--220. The proof under the generalised Riemann hypothesis.
\item W.~Diffie and M.~E.~Hellman, New directions in cryptography, \emph{IEEE Transactions on Information Theory} 22 (1976), 644--654. The paper that turned discrete logarithms into public-key cryptography.
\item P.~Moree, \emph{Artin's primitive root conjecture -- a survey}, Integers 12 (2012), A13. The 1927 conversation with Hasse, Artin's constant $0.3739\ldots$, and Heath-Brown's theorem (Quart.\ J.\ Math.\ 37 (1986), 27--38).
\end{itemize}
